\section{Introduction}

Protein-to-genome alignment---the task of mapping a protein sequence to its encoding genomic locus and producing a spliced gene structure---is essential for genome annotation, orthology inference, and evolutionary analysis \cite{birney2004gene,stanke2005augustus}. As genome sequencing initiatives scale to encompass the tree of life~\cite{lewin2018ebp}, automated annotation pipelines depend critically on accurate protein-to-genome aligners as a primary evidence source~\cite{gabriel2024braker3,bruna2023galba,keilwagen2018gemoma}. Over two decades of tool development have produced mature aligners spanning diverse algorithmic strategies: minimizer-based spliced alignment (miniprot~\cite{li2023miniprot}), profile-HMM methods (Spaln3~\cite{gotoh2024spaln3}), and dynamic programming approaches (Exonerate~\cite{slater2005exonerate}, GeneWise~\cite{birney2004genewise}), among others. On the standard benchmark of one-to-one ortholog alignment, these tools perform well: the best tools achieve exon F1 $>$ 0.87 across mammals (Table~\ref{tab:benchmark}). By that measure, one-to-one ortholog alignment is mature.

However, most protein-coding genes belong to paralogous families~\cite{ohno2013evolution}, arising from the two rounds of whole-genome duplication in early vertebrate evolution~\cite{dehal2005tworounds}, where one protein matches multiple genomic loci in a target species. In these one-to-many cases---the biologically common scenario---we find that paralog-rich regions represent a high-risk zone for existing tools. Across six tools and four vertebrate species, primary exon F1 drops to 0.17--0.40 on paralogs, and the correct-locus rate collapses from 0.91--0.99 (one-to-one) to 0.11--0.49 (paralog). The performance gap between one-to-one and paralog regimes is 55--80\% across all tools and species (Table~\ref{tab:benchmark}). This is not a tool-specific weakness; it is a systematic vulnerability of single-tool, single-query approaches in regions with multiple homologous copies.

Why does this happen? Prior benchmarking studies~\cite{burset1996evaluation} evaluated gene prediction tools on single-gene regions, not genome-scale protein queries with paralog ambiguity. The standard evaluation uses one-to-one orthologs between closely related species (e.g., human--mouse), where both high sequence identity and locus uniqueness make the task easy. Systematic evaluation on paralogs---and more importantly, causal diagnosis of \emph{why} paralogs fail---has been absent from the literature.

We conduct this diagnosis through a forced-locus experiment that decouples search error from alignment error. The results decompose the paralog failure into two independent, additive root causes:

\begin{enumerate}
\item \textbf{Cross-species identity ceiling (RC1):} Even when given the correct genomic locus, a human protein query against a non-human paralog achieves only 0.64--0.71 exon F1 with miniprot, compared to 0.92--0.98 for forced-locus one-to-one orthologs (Table~\ref{tab:forced_locus}). The 0.22--0.30 gap is an intrinsic limit imposed by amino acid divergence: the human query protein lacks the sequence motifs needed to resolve exon boundaries in a diverged target species' paralog. No C-stage alignment algorithm can overcome this bottleneck---the query itself is the problem.

\item \textbf{Locus confusion (RC2):} In the paralog regime, the correct-locus rate collapses from 0.91--0.99 to 0.11--0.49 across tools and species (Fig.~\ref{fig:forced_locus}b). Tools select the paralog copy with the highest raw sequence similarity to the query, which is not necessarily the evolutionarily correct target. This accounts for an additional 0.22--0.38 F1 loss beyond the identity ceiling.
\end{enumerate}

These findings directly motivate the design of TPRS. RC1 implies that the candidate search space must include query proteins with high sequence identity to the target---specifically, the target species' own protein sequence. RC2 implies that (a) the candidate space must include genomic regions found by diverse search strategies (no single tool finds all correct loci), and (b) the selection among candidates must be learned rather than heuristic, because the discriminating signal is distributed across multiple features with non-linear interactions.

The diagnosis leads to a clear design principle: \textbf{diagnose the ceiling of annotation-free approaches, then introduce annotation-derived priors that break through it}. TPRS is an alignment method that follows this principle in two stages. \textbf{Stage 1---Candidate Generation}---expands the search space along three orthogonal dimensions: tool diversity (4 complementary A-stage tools with chromosome- and strand-aware merging), query identity (both human and species-matched protein queries against each candidate locus), and prior knowledge (Ensembl Compara co-ortholog recovery for proteins missed by all tools). This produces a candidate pool of (locus, query-type) pairs per protein. \textbf{Stage 2---Learned Joint Selection}---scores every candidate pair using an XGBoost regressor trained on 25 features across 5 categories, and selects the pair with the highest predicted exon F1. The resulting spliced alignment is the direct output—TPRS competes on alignment metrics, not annotation metrics, and its design is the logical consequence of the causal diagnosis: where sequence evidence alone is provably insufficient, annotation-derived priors make the problem tractable.

Our contributions are threefold:
\begin{itemize}
\item \textbf{Causal diagnosis:} We provide the first systematic decomposition of paralog alignment failure into two independent root causes, with quantified F1 loss attribution across 4 species and 6--7 tools.
\item \textbf{Alignment method:} We present TPRS, a diagnosis-driven alignment method (Candidate Generation + Learned Joint Selection) whose architecture follows directly from the causal analysis. Each component addresses a specific diagnosed failure mode, and each is independently validated through ablation. TPRS demonstrates that annotation-derived auxiliary information---species-matched proteins and curated orthology---can be used to transform an alignment problem from broken to usable, without requiring a full annotation pipeline.
\item \textbf{Empirical validation:} Within species with existing annotation resources, TPRS achieves paralog primary exon F1 of 0.57--0.79, improving $+$0.36--0.46 over the best single tool, and recovers 66--90\% of the one-to-one ceiling. All 32 ablation outcomes share the same sign across species ($p < 5\times10^{-10}$), establishing that the design generalizes.
\end{itemize}
